\section{Il libro scritto da Augusto Righi}

L'ottica delle oscillazioni elettriche \cite{righi1897ottica} è un trattato monografico scritto da Augusto Righi nel 1897, in cui vengono descritte le 
esperienze da lui condotte per verificare le proprietà delle onde elettromagnetiche.\\
Mario Gliozzi (storico della fisica) lo descrive come il "testo che chiude un'epoca", in quanto raccoglie in un unico volume 
tutte le esperienze fatte da Hertz e ne verifica le analogie con le esperienze di ottica classica.\\
Il volume è diviso in tre parti:

\begin{itemize}
    \item la prima parte (da pagina 3 a pagina 71) parla della costruzione degli oscillatori e dei risonatori ideati da Righi, a partire dalle esperienze di Hertz \cite{bryant1988heinrich}
    \item la seconda parte (da pagina 75 a pagina 194), divisa in 7 capitoli, descrive le varie esperienze con onde emesse dai propri emettitori, ideate in analogia con l'ottica (interferenza, diffrazione, trasparenza/Opacità, riflessione, rifrazione/riflessione totale, doppia rifrazione)
    \item le note (da pagina 197 a pagina 248) costituiscono la parte finale del volume e contengono tutte le formule matematiche utilizzate da Righi nelle esperienze, ma analizzate più in dettaglio.
\end{itemize}    







